\chapter[Extracellular fields and inverse problems]{From membrane currents to extracellular electrodes}
\label{ch:electrode-physics}
An extracellular electrode does not measure a neuron's voltage
in isolation. It measures a difference of electric potential
created by spatially distributed ionic currents, tissue geometry,
conductivity, referencing, and acquisition electronics
\cite{buzsaki2012,einevoll2013}.

\section{Electrochemical driving forces}
For an ideal dilute permeant ion with valence $z\ne0$, equality
of electrochemical potentials on the two sides of a membrane gives
\[
E_{\mathrm{ion}}=\frac{RT}{zF}
  \log\frac{c_{\mathrm{out}}}{c_{\mathrm{in}}}.
\]
This is the Nernst equilibrium potential under its concentration
and temperature assumptions. The Goldman--Hodgkin--Katz formula
extends the constant-field model to multiple monovalent ionic
species with permeability parameters; it is not equivalent to
measuring those permeabilities from a single LFP.
In the Hodgkin--Huxley description the sodium current is modeled
as $I_{\mathrm{Na}}=\bar g_{\mathrm{Na}}m^3h(V-E_{\mathrm{Na}})$;
gating variables obey voltage-dependent kinetic equations.
The passive cable approximation has the schematic form
\[
c_m\partial_t V(x,t)=g_a\partial_{xx}V(x,t)
   -g_m(V-E_L)+i_{\mathrm{ext}}(x,t),
\]
with coefficients per unit length and consistent units.
Cable and membrane models determine transmembrane sources; they
are not by themselves electrode forward models.

\section{Quasistatic field equation}
At electrophysiological frequencies and spatial scales where
inductive effects are negligible, write $\mathbf E=-\nabla\phi$.
Let impressed current $\mathbf J_p$ and Ohmic return current
$-\boldsymbol\sigma\nabla\phi$ sum to a divergence-free total
current. Current conservation implies
\[
\nabla\!\cdot(\boldsymbol\sigma\nabla\phi)
       =\nabla\!\cdot\mathbf J_p,
\quad\text{or}\quad
-\nabla\!\cdot(\boldsymbol\sigma\nabla\phi)=q,
\quad q=-\nabla\!\cdot\mathbf J_p.
\]
The tensor $\boldsymbol\sigma$ may be anisotropic and spatially
inhomogeneous. To solve the equation, the domain, conductivity,
boundary conditions and electrode reference must be specified.
For an infinite isotropic medium, a point source of strength $I$
located at $\mathbf r_0$ gives, away from that point,
\[
\phi(\mathbf r)=\frac{I}{4\pi\sigma
  |\mathbf r-\mathbf r_0|}.
\]
Superposition yields a discrete source model $\mathbf v=G\mathbf s$
when the forward physics is linear, with lead-field entries obtained
by integrating the Green kernel over chosen source basis functions.
A dipole arises as the leading far-field contribution when net
source current cancels; line-source models integrate along a dendrite.

\begin{proposition}[Uniqueness for fixed Dirichlet conditions]
Let $\Omega$ be a bounded connected region and
$\boldsymbol\sigma$ uniformly positive definite. If two weak
solutions of the same elliptic equation have equal Dirichlet
boundary values, they agree (almost everywhere).
\end{proposition}
\begin{proof}
Their difference $w$ satisfies the homogeneous equation and
zero boundary values. Multiply by $w$, integrate by parts,
and use the boundary condition:
$\int_\Omega\nabla w\cdot\boldsymbol\sigma\nabla w\,dx=0$.
Positivity implies $\nabla w=0$; the boundary forces $w=0$.
For pure Neumann conditions, uniqueness holds only up to an
additive constant, with a compatibility condition on source flux.
\end{proof}

\section{Why the inverse need not be unique}
Suppose an electrode array measures $\mathbf v=G\mathbf s$.
If $G\mathbf h=0$ for nonzero $\mathbf h$, then
$G(\mathbf s+\mathbf h)=G\mathbf s$. This is exact
finite-dimensional nonidentifiability. In the continuum, smoothing
operators can also be ill-conditioned or compact, so tiny
measurement errors may cause enormous source errors even when
uniqueness holds in a restricted model.
\emph{Hadamard well-posedness} requires existence, uniqueness
and continuous dependence of solutions on data. A linear operator
is \emph{compact} if it maps bounded sets into sets with compact
closure; its singular values in infinite-dimensional inverse problems
can converge to zero.
For singular system $Gv_j=\sigma_j u_j$ and data $b$,
the formal inverse contains $\langle b,u_j\rangle/\sigma_j$
and magnifies noise when
$\sigma_j\to0$. A classical Picard-type condition controls the
squared normalized coefficients for compatible data.
Tikhonov regularization instead solves
\[
\widehat{\mathbf s}_\lambda
=\arg\min_{\mathbf s}
\{\|G\mathbf s-\mathbf v\|_2^2
 +\lambda\|\mathbf s\|_2^2\}
=(G^\top G+\lambda I)^{-1}G^\top\mathbf v.
\]
\begin{example}[MATLAB and interpretation]\sloppy
\path{field_potential_forward.m} receives $G$;
it does \emph{not} compute tissue conductivity or the Green kernel.
\path{field_potential_inverse_tikhonov.m}
returns one regularized source estimate, not a unique membrane-current
pattern. \path{current_source_density.m} needs
spatial geometry and differentiated potentials; references and
volume conduction remain crucial.
\end{example}
% --- expanded edition (October 2026) ---
\section{The Goldman--Hodgkin--Katz equations}
The constant-field model assumes that the electric field is uniform
across a membrane of thickness $d$, that ions cross independently, and
that the flux of each species obeys the Nernst--Planck equation
$J=-D(\partial_xc+\frac{zF}{RT}c\,\partial_x\phi)$. With
$u=zFV/(RT)$ and permeability $P=D\beta/d$ ($\beta$ a partition
coefficient), integrating across the membrane gives the
\emph{GHK current equation}
\[
I=P z^2\frac{F^2V}{RT}\;\frac{c_{\rm in}-c_{\rm out}e^{-u}}{1-e^{-u}}.
\]
The current vanishes at $c_{\rm in}=c_{\rm out}e^{-u}$, i.e.\ at the
Nernst potential, and is a nonlinear (rectifying) function of $V$ when
concentrations differ. Setting the \emph{total} current of K$^+$,
Na$^+$ and Cl$^-$ to zero gives the \emph{GHK voltage equation}
\[
V_{\rm rest}=\frac{RT}{F}\log\frac{P_{\rm K}[K]_{\rm o}+P_{\rm Na}[Na]_{\rm o}+P_{\rm Cl}[Cl]_{\rm i}}{P_{\rm K}[K]_{\rm i}+P_{\rm Na}[Na]_{\rm i}+P_{\rm Cl}[Cl]_{\rm o}}.
\]
\begin{example}[Resting potential]
With $[K]_{\rm o}=4$, $[K]_{\rm i}=140$, $[Na]_{\rm o}=145$,
$[Na]_{\rm i}=12$, $[Cl]_{\rm o}=120$, $[Cl]_{\rm i}=10$\,mM, relative
permeabilities $1:0.04:0.45$ and $RT/F\approx26.7$\,mV at
$37^\circ$C, the numerator is $4+5.8+4.5=14.3$ and the denominator
$140+0.48+54=194.5$, giving $V_{\rm rest}\approx26.7\ln(0.0735)\approx-70$\,mV.
The potassium Nernst potential is $26.7\ln(4/140)\approx-95$\,mV. The
rest is depolarized relative to $E_{\rm K}$ because of the sodium and
chloride terms.
\end{example}
The GHK equations are a model, not a law: they fail when channels are
multiply occupied, when ions interact, or when the field is not uniform.
Their ratio form also shows that only \emph{relative} permeabilities
are determined by a reversal potential, so a single measured reversal
cannot fix three permeabilities.

\section{The Hodgkin--Huxley system}
The complete space-clamped model is
\[
\begin{aligned}
C_m\dot V&=-\bar g_{\rm Na}m^3h(V-E_{\rm Na})-\bar g_{\rm K}n^4(V-E_{\rm K})-g_L(V-E_L)+I_{\rm ext},\\
\dot x&=\alpha_x(V)(1-x)-\beta_x(V)x,\qquad x\in\{m,h,n\}.
\end{aligned}
\]
Each gate relaxes to $x_\infty(V)=\alpha_x/(\alpha_x+\beta_x)$ with time
constant $\tau_x(V)=1/(\alpha_x+\beta_x)$. Sodium activation $m$ is fast,
inactivation $h$ and potassium activation $n$ are slow. The spike arises
because $m$ jumps quickly after a depolarization, giving positive
feedback through $E_{\rm Na}$, while the slower decrease of $h$ and
increase of $n$ terminate it. Reducing the model by setting
$m=m_\infty(V)$ and exploiting the approximate relation $h+n\approx0.8$
yields a two-dimensional system whose phase plane resembles
FitzHugh--Nagumo; repetitive firing begins at a Hopf bifurcation of
its equilibrium (Chapter~\ref{ch:oscillator-foundations}).

For extracellular recording the essential output is the sum of
transmembrane currents per unit area,
$i_m=C_m\dot V+i_{\rm ion}$. By charge conservation along a neuron, the
integral of $i_m$ over the whole cell surface is zero at every instant:
a neuron has no net monopole source, and its far field is at least
dipolar.

\section{Cable equation: space and time constants}
Write $r_m$ for membrane resistance times length (units $\Omega\,\mathrm m$),
$c_m$ for membrane capacitance per length ($\mathrm F/\mathrm m$), and
$r_a$ for axial resistance per length ($\Omega/\mathrm m$). Thus a segment
of length $dx$ has membrane resistance $r_m/dx$, capacitance $c_m dx$,
and axial resistance $r_a dx$. The passive cable equation is
$\lambda^2\partial_{xx}V-\tau\partial_tV-V=-r_m i_{\rm ext}$,
where $i_{\rm ext}$ is injected current per length,
$\lambda=\sqrt{r_m/r_a}$ has units of length and $\tau=r_mc_m$ has units
of time. For a point injection at $x=0$ in an infinite passive cable,
the steady solution is $V(x)=V(0)e^{-|x|/\lambda}$.

For the convention $V(x,t)=\Re\{\widehat V(x)e^{i\omega t}\}$, away
from the injection the equation gives
\[
\widehat V''=\kappa^2\widehat V,\qquad
\kappa=\frac{\sqrt{1+i\omega\tau}}{\lambda},\quad\Re\kappa>0,
\qquad \widehat V(x)=\widehat V(0)e^{-\kappa|x|}.
\]
The complex length is $1/\kappa$; the real amplitude attenuation length is
\[
\lambda_\omega=\frac1{\Re\kappa}
=\lambda\sqrt{\frac{2}{1+\sqrt{1+\omega^2\tau^2}}}.
\]
The imaginary part of $\kappa$ gives spatial phase lag. At zero frequency
$\lambda_\omega=\lambda$; for $|\omega\tau|\gg1$,
$\lambda_\omega\sim\lambda\sqrt{2/|\omega\tau|}$.
These conclusions assume a passive, uniform infinite cable; branches,
finite boundary conditions and active channels change the transfer function.
The MATLAB derivation regression compares this solution to a converging
finite-difference cable calculation, not to a fitted biological neuron.

\section{The Picard condition and regularization parameter choice}
Let the discretized forward operator have SVD $G=\sum_j\sigma_ju_jv_j^\top$
and data $\mathbf v=G\mathbf s_{\rm true}+\boldsymbol\epsilon$. The
\emph{discrete Picard condition} requires that the coefficients
$|u_j^\top\mathbf v|$ decay, on average, faster than $\sigma_j$ until they
reach the noise level. Plotting $\sigma_j$, $|u_j^\top\mathbf v|$ and
their ratio on a logarithmic scale (a \emph{Picard plot}) shows which
components are informative: beyond the index where
$|u_j^\top\mathbf v|$ levels off at the noise floor, the ratio grows
and those components should be filtered. The Tikhonov estimate has
filter factors $\sigma_j^2/(\sigma_j^2+\lambda)$; common choices of
$\lambda$ are:
\begin{itemize}[leftmargin=2em]
\item the \emph{discrepancy principle}: choose $\lambda$ so that
$\|G\hat{\mathbf s}_\lambda-\mathbf v\|\approx\|\boldsymbol\epsilon\|$,
which requires an estimate of the noise level;
\item the \emph{L-curve}: plot $\log\|\hat{\mathbf s}_\lambda\|$ against
$\log\|G\hat{\mathbf s}_\lambda-\mathbf v\|$ and choose the corner;
\item \emph{generalized cross-validation}, which minimizes
$\|G\hat{\mathbf s}_\lambda-\mathbf v\|^2/\tr(I-A_\lambda)^2$ with
influence matrix $A_\lambda=G(G^\top G+\lambda I)^{-1}G^\top$.
\end{itemize}
Each criterion is a statistical decision with known failure modes: the
L-curve can lack a corner, GCV can undersmooth with correlated noise,
and the discrepancy principle depends on a noise estimate.

\section{Line sources and laminar forward models}
Assume an infinite homogeneous isotropic conductor with conductivity
$\sigma>0$. A straight segment from $\mathbf r_1$ to $\mathbf r_2$ has
length $\ell>0$ and direction $\mathbf u=(\mathbf r_2-\mathbf r_1)/\ell$.
Let $z=(\mathbf r-\mathbf r_1)\cdot\mathbf u$ and
$\rho=\|\mathbf r-\mathbf r_1-z\mathbf u\|>0$.
For total current $I$ distributed uniformly along the segment,
positive outward into the conductor, integrating the point-source kernel gives
\[
\begin{aligned}
\phi(\mathbf r)
&=\frac{I}{4\pi\sigma\ell}\int_0^\ell
 \frac{ds}{\sqrt{\rho^2+(z-s)^2}}\\
&=\frac{I}{4\pi\sigma\ell}
 \left[\operatorname{asinh}\frac{\ell-z}{\rho}
       +\operatorname{asinh}\frac{z}{\rho}\right].
\end{aligned}
\]
With $h=-z$ and $l=\ell-z=h+\ell$, this equals
$\frac{I}{4\pi\sigma\ell}
\log\frac{\sqrt{l^2+\rho^2}+l}{\sqrt{h^2+\rho^2}+h}$.
All distances are in metres, $I$ in amperes and $\sigma$ in siemens per
metre, yielding volts. Reversing the segment endpoints leaves the potential
unchanged. Holding total current fixed and shrinking a centred segment
recovers $I/(4\pi\sigma r)$ away from the source.

The ideal zero-radius line still has a logarithmic singularity as
$\rho\to0$ inside the segment; finite-radius or electrode-averaging models
are needed there. It does not eliminate all near-source singularities.
At distances large compared with $\ell$, the isolated segment behaves as
a point source; dipolar cancellation requires the return-current geometry.
The supplied MATLAB forward operator only multiplies a user-provided
lead field; it does not construct this geometry automatically.

\section*{Exercises for this chapter}
\begin{exercise}
For equal concentrations $c_{\rm in}=c_{\rm out}=c$, simplify the
GHK current equation, taking the removable limit at $V=0$, and show
that $I=Pz^2F^2cV/(RT)$. Explain why unequal concentrations
usually yield a voltage-dependent effective conductance and why a
local linearization about a reversal potential is not a global Ohmic law.
\end{exercise}
\begin{exercise}
Compute the potassium, sodium and chloride Nernst potentials for the
example concentrations.
\end{exercise}
\begin{exercise}
For $\tau=20$\,ms, compute $\lambda_\omega/\lambda$ at $1$, $10$ and
$100$\,Hz.
\end{exercise}
\begin{exercise}
Derive the line-source formula by integrating
$I\,ds/(4\pi\sigma\ell|\mathbf r-\mathbf r(s)|)$ along the segment.
\end{exercise}
